\vspace{-5pt}
\section{Conclusion and Discussion}
\vspace{-7pt}
\label{sec:discussion}
% In this work, we introduce a novel downstream performance scaling framework including (1) the COD method that effectively models the underlying diverse scaling patterns of evaluation tasks
\camera{In this work, we propose a novel framework for predicting LLM downstream performance scaling, which consists of three key contributions: (1) the COD framework that effectively models the intrinsic diverse scaling patterns of tasks in the evaluation set}; (2) a scaling law for downstream task performance that provides a fitting formula for performance-compute extrapolation; and (3) a systematic methodology for identifying and leveraging predictable subset that provides a robust intermediate metric for accurate full-set performance predictions.
We discuss the limitations and future works in Appendix~\ref{sec:limiations}.
% Looking forward, our approach can be further expanded across model architectures, training methods, and evaluation set types, while extending this framework to address chain-of-thought reasoning patterns offer promising avenues for future research.

% \paragraph{Boarder Impact} 
% \section{Boarder Impact}
% The LLMs performance scaling prediction method proposed in this work can be used to guide decision-making for training LLMs and to monitor whether the LLM training process aligns with expectations. Accurate prediction allows for better planning of compute and human resources, thereby playing a valuable role in scientifically advancing the training of larger LLMs. However, if the prediction results are inaccurate, it may lead to outcomes that do not meet expectations for the invested resources, potentially resulting in resource waste or economic losses.

\section*{Ethics Statement}
\label{sec:social_impact}
We have read and adhered to the ICLR Code of Ethics. This work proposes a computational framework to enable more efficient resource allocation in the training of LLMs. The research is methodological in nature and aims to support more sustainable and responsible practices within the field of AI development.

We provide a detailed account of our methods, theoretical proofs, and experimental settings. We openly discuss the limitations of our framework in Appendix~\ref{sec:limiations}. This study does not involve human subjects or the use of sensitive personal data.

We utilized language models at the writing level, including checking for grammatical errors in the article and modifying expressions. The use of language models had no impact on the article's innovative contributions, experiments, or analytical perspectives.



\section*{Reproducibility Statement}
We have made extensive efforts to ensure the reproducibility of our work. The core methodology, the Clustering-On-Difficulty (COD) framework, is detailed in \cref{sec:method}. The improved MeanShift clustering algorithm is described in \cref{sec:method-clustering} , with full pseudocode provided in Appendix~\ref{sec:improved} (Algorithm 1). Our performance scaling law (Theorem 1) is presented in \cref{sec:fitting} , with a complete proof available in \cref{sec:prop_proof}. We discuss the additional computational cost of the COD method in Appendix~\ref{append:compute}.

Our experimental setup, including model architectures, training data philosophy, and hyperparameters, is thoroughly documented in \cref{sec:exp_setups} and Appendix~\ref{append:setting} , with specific model configurations listed in \cref{tab:model_configs_dense}. The evaluation benchmarks, protocols, and few-shot settings are described in \cref{sec:exp_setups} and summarized in \cref{tab:evaluation_info}. Extensive ablation studies validating our component choices—including extrapolation formulas (\cref{sec:ablation_extrapolation}), clustering algorithms (Appendix~\ref{sec:clustering_ablation}), interpolation methods (Appendix~\ref{sec:interpolation_ablation}), and criteria for filtering clusters (Appendix~\ref{sec:criteria})—are provided to support our findings. We visualize the task difficulty distribution for each evalset in Appendix~\ref{sec:difficulty_distribution}.